\chapter{Follow the audio across its boundaries}
\MARSChapterHeading{Follow the audio across its boundaries}
\label{revision-lab:wav}

Read Chapter~23 alongside this lab. Import the native lesson
\emph{Follow audio across its boundaries}, built from
\path{revision/lessons/031-pcm-wav}. Its \texttt{worked} directory contains
five complete experiments. Predict the result inline, then use ordinary Debug
to inspect the named boundaries. Each label identifies the next instruction:
read the state before executing that instruction. Finish the process before
starting another activity. Figure~\ref{fig:revision-wav-native-prediction}
shows an inline question and its source pointer.

The continuing source is \texttt{independent/wav-reverse-frames.asm} in the
imported lesson. Its caller and six WAV helpers are complete; only
\texttt{wav\_reverse\_frames} is unfinished. This source includes the supplied
helper directly, so it needs no preparation command or neighboring include.
The three assessment groups cover twelve mono cases, twelve stereo cases and
eleven boundary, reuse and pipeline cases. Recheck all 35 after your final edit.
Graduated hints and free-writing mode follow your instructor's release policy.

The split-file worked program remains at
\path{revision/labs/files/wav-pipeline.asm}, with
\texttt{wav-functions.asm} beside it. The split-file independent source is
\texttt{wav-reverse-frames.asm} in that same directory. Commands below prepare
additional standalone observations of those files; the native lesson's worked
experiments already contain their observation data. The original six-function
construction task remains in \texttt{packages/031-pcm-wav/} if you want to
implement the supplied helpers as a separate exercise.

\begin{figure}[!ht]
\centering
\begin{tikzpicture}[alt={Two crops from one actual light-theme IDE photograph. The inline quiz asks where the next header begins for a chunk at byte twelve with an eight-byte header, three payload bytes and one pad byte. Choices are 24, 23 and 20. Label A marks the inline answer area. The source crop shows the padded extent calculation and a yellow lesson pointer at the next full-width length experiment. Label B identifies that source pointer. No program is running in this photograph.}]
\node[anchor=north west,inner sep=0] (quiz) at (0,0) {\MARSRevisionCrop{revision/captures/wav/wav-native-prediction.png}{0.94\linewidth}{0.04843750}{0.37333333}{0.59609375}{0.38444444}};
\begin{scope}[shift={(quiz.south west)},x={($(quiz.south east)-(quiz.south west)$)},y={($(quiz.north west)-(quiz.south west)$)}]
\draw[black,line width=.7pt] (.065,.06) rectangle (.95,.53);
\node[draw=black,fill=white,inner sep=1pt,font=\small\bfseries] at (.025,.30) {A};
\end{scope}
\node[anchor=north west,inner sep=0] (code) at ([yshift=-3mm]quiz.south west) {\MARSRevisionCrop{revision/captures/wav/wav-native-prediction.png}{0.91\linewidth}{0.42968750}{0.58888889}{0.22656250}{0.21333333}};
\begin{scope}[shift={(code.south west)},x={($(code.south east)-(code.south west)$)},y={($(code.north west)-(code.south west)$)}]
\draw[black,line width=.7pt] (.065,.01) rectangle (.57,.145);
\node[draw=black,fill=white,inner sep=1pt,font=\small\bfseries] at (.025,.075) {B};
\end{scope}
\end{tikzpicture}
\caption{Answer the prediction directly in the Learn pane (A). The next header
begins at $12+8+3+1=24$. The source has already expressed that calculation;
its yellow pointer (B) identifies the next instruction to inspect in this
experiment. A lesson pointer is a source reference, not evidence that the
instruction has executed. No Debug process is running in this photograph.
Both crops retain the original capture pixels; the outlines and labels are
native LaTeX annotations.}
\label{fig:revision-wav-native-prediction}
\end{figure}


\section{Predict the parser's next address}

Generate the trace using the main chapter's fixture command. On paper, begin
at chunk offset twelve and calculate each successive header using the declared
payload length and odd-byte padding. Keep the complete 82-byte extent separate
from the sixteen-byte PCM extent. Explain why the first sample begins at 56,
why a data length of sixteen means four stereo frames, and why the parser must
continue at offset 72 after finding those samples.

\Needspace{8\baselineskip}
In the native lesson, Debug \texttt{worked/parser-publication.asm}. At
\texttt{parsed\_record}, \texttt{s0} through \texttt{s4} contain the five
fields below. Continue to \texttt{rejected\_record}: the second parse has changed
the trailing chunk's declared length from one to three without extending the
file. All four eight-byte sentinel groups in \texttt{s0} through \texttt{s3}
remain \texttt{a5}; status is one and returned count is zero. Finding a valid
data chunk earlier does not permit publishing a partial record.

\Needspace{10\baselineskip}
For the split-file observation, prepare step \texttt{s06-1}:
\Needspace{8\baselineskip}
\MARSCodeLabel{Select a parser observation without changing instructions}
\begin{Verbatim}
python3 packages/031-pcm-wav/prepare.py s06-1 \
  --source revision/labs/files/wav-pipeline.asm \
  --output revision/labs/files/wav-parse.asm
\end{Verbatim}
Set a breakpoint at \texttt{parse\_return} and Debug the prepared caller.
Before continuing, interpret the metadata bytes as the five fields in the
main chapter. The expected record is $(16000,2,4,56,16)$, with widths
$(4,4,8,8,8)$ bytes. A display of \texttt{38} at descriptor offset sixteen
means a file-relative displacement, not a pointer to address hexadecimal 38.
The parser returns four frames in \texttt{v1} at this boundary.

Inspect a second observation whose trailing \texttt{note} claims three payload
bytes while the complete count remains 82. Predict the required next offset
before running. Initialize metadata to a recognizable pattern, then check that
all 32 bytes retain it after format error one. If fields change, locate the
first store to the public record and ask what validation still remains there.
Merely obtaining an error status does not satisfy the publication guarantee.

For an additional boundary experiment, let an unknown chunk declare length
hexadecimal \texttt{ffffffff}. Derive its padded extent in 64-bit arithmetic.
State which comparison must reject it before any payload access. This is a
parser arithmetic test, not a request to allocate a multi-gigabyte input.

\section{Separate the two channel histories}

The native experiment \texttt{worked/channel-histories.asm} exposes all eight
signed output samples in \texttt{s0} through \texttt{s7} at
\texttt{filtered\_frames}. Pair adjacent registers as left and right; the
second frame is $(1500,500)$, not $(0,2000)$, because the first raw samples
remain in their respective histories. At \texttt{output\_guard}, the next
byte is still hexadecimal \texttt{a5}. This one sentinel is useful evidence;
it does not by itself prove that every earlier access was permitted.

For the split-file observation, prepare step \texttt{s29-1} and pause at
\texttt{filter\_return}. Decode each
pair of bytes as one signed sample, then group adjacent pairs into stereo
frames. Derive the first two frames by hand using seven zeros of history per
channel. A left output influenced by the right input indicates that the gather
address has confused a frame stride with a sample stride.

For a more discriminating experiment, use unequal coefficients and a left-only
impulse. An all-equal average cannot expose reversed coefficient order, and
identical channels cannot expose accidental cross-channel mixing. Predict
which coefficient meets the impulse at the first output, then follow that
sample as the window moves.

Break at \texttt{filter\_history} after a block. Identify the $7+m$ input
samples in the temporary window and select its final seven. For $m=3$, write
which four come from the old history and which three are new. Then inspect a
recording longer than 64 frames on both sides of that block boundary. If the
second block begins with another zero-history transient, the caller has reset
state too early. If it uses preceding filtered values, it has introduced a
recurrence that the FIR contract never requested.

\section{Inspect the file before saving it}

Debug \texttt{worked/encoder-capacity.asm} in the native lesson. Its first
call offers 59 bytes for a 60-byte image. At \texttt{capacity\_rejected},
status two accompanies an unchanged first destination group. The next call
offers 60 bytes and succeeds. At \texttt{encoded\_image}, distinguish
RIFF size 52, byte rate 64000, four bytes per frame and sixteen PCM bytes.
The first output sample is 1500 and the byte beyond the complete image remains
\texttt{a5}. The supplied input here is the filtered trace.

The original split-file encode observation is step \texttt{s43-1}. It selects PCM already
stored in the caller's \texttt{samples} arena; an encode-only run does not first
filter it. To inspect the filtered trace's complete output, use pipeline step
\texttt{s63-1} with the generated input and pause at \texttt{encode\_return}.
This distinction matters when comparing an isolated function with a composed
program.

Locate the four-byte data length and the first PCM byte at output offset 44.
For four stereo frames, predict the complete byte count, RIFF size, block
alignment and byte rate before reading them. Check the first bytes beyond the
sixty-byte image as well. A correct header does not excuse writing an extra
sample into the guard region.

Try a capacity of 59 for this image and a destination filled with sentinels.
The encoder must report capacity error two without changing any selected
caller memory. A valid zero-frame prefix instead needs 44 bytes and no PCM
access. Neither case requires a special fake sample to make the format work.

Now trace the load and save loops. At each successful transfer, identify the
remaining count and next byte pointer. Explain why a positive short read is
not yet EOF, why a full input buffer needs a separate one-byte probe, and why
a zero-byte write cannot be retried indefinitely. Test malformed input against
an existing output file: validation should fail before output opening. Keep
that guarantee separate from a later disk-write error, which may leave a
partially written file.

\section{Write a frame reversal that can share its buffer}

Implement \texttt{wav\_reverse\_frames} in the independent caller. Its arguments
are PCM pointer in \texttt{a0}, eight-byte-aligned descriptor prefix pointer in
\texttt{a1}, and output PCM pointer in \texttt{a2}. The prefix contains the
same rate, channel and full-width frame count as the worked program; validate
it under that program's profile. The selected PCM length is $2CN$. Pointers
may be byte-unaligned. Input and output are either disjoint or exactly equal;
other overlap is outside the contract. The prefix is disjoint from both.

On success, return zero in \texttt{v0} and original frame count $N$ in
\texttt{v1}. Output frame $i$ must equal original frame $N-1-i$, with all
$2C$ bytes within each frame unchanged. On profile failure, return one and zero
without changing caller objects. For a valid zero-frame prefix, return success
without touching either PCM pointer. You may call \texttt{wav\_bound} to
validate the prefix, but its successful byte-count result is not this function's
frame-count result.

Preserve every bit of \texttt{s0} through \texttt{s7}, \texttt{fp}, \texttt{sp}
and the supplied return continuation. The incoming stack is eight-byte aligned;
restore it on each exit. Write only the selected output interval and your
private stack frame. Disjoint input remains unchanged; exact aliasing permits
changes only within that common PCM interval. Do no console or file I/O in the
new function. Caller-saved scalar and vector registers need not be retained.

Begin with three mono frames whose two-byte codes are all distinct, then three
stereo frames whose left and right values differ. Write down the expected bytes.
Reversing individual bytes would exchange each sample's significance. Reversing
individual samples would exchange the channels within stereo frames. Both can
sound different while still producing a superficially valid WAV header.

For disjoint output, one possible loop reads frames from the end and writes
from the beginning. For exact aliasing, reason about an endpoint pair before
storing either one. You can exchange a frame byte at a time with two loaded
bytes, or exchange its channels as complete two-byte codes. The odd middle
frame needs no change. Establish a loop invariant that states which endpoint
frames are final and which original frames are still available.

\Needspace{7\baselineskip}
The provided caller invokes your function after successful filtering, using
\texttt{samples} as both source and destination, then encodes and saves.
For the trace, the sixteen output PCM bytes must be
\texttt{f4 11 f4 f1 f4 01 f4 01 dc 05 f4 01 dc 05 24 fa}.
The output still has four stereo frames at 16000 Hz. Its prefix needs no change.
Its first frame comes from the end of the filtered sequence, not from the end
of the unfiltered input.

\section{Check the function and its integration separately}

Use the inline \emph{Check this program} controls for the native lesson's
continuing source. Its full-pipeline case supplies virtual input and output,
limits reads to seven bytes and writes to five, and begins descriptor allocation
at 19. Your reversal performs no I/O; the supplied load and save helpers own
those operations. The expected saved file is exactly sixty bytes. A successful
reversal must also work in the complete program.

An ordinary Run of the independent source selects its default bound calculation
and terminates silently. It does not exercise reversal or save a WAV file.
Use the full native assessment for the integrated check, then the generated
host-file experiment below as an additional test. After implementation edits,
use the final full recheck to renew results for that same source.

For command-line work, the public behavioral check is
\path{revision/labs/checks/wav-reverse-frames.kts}. It selects different valid
prefixes, byte offsets, exact-alias and disjoint buffers, invalid full-width
counts, repeated calls and a complete file pipeline. It compares state and
bytes, without requiring the loop above or particular temporary registers.
Declare the helper explicitly when using the grader, which freezes only the
sources you supply:
\Needspace{9\baselineskip}
\MARSCodeLabel{Grade the independent caller with its include dependency}
\begin{Verbatim}
java -jar "$CLI" grade \
  revision/labs/checks/wav-reverse-frames.kts \
  revision/labs/files/wav-reverse-frames.asm \
  --include revision/labs/files/wav-functions.asm
\end{Verbatim}
A deliberate \texttt{break} marks the unfinished entry. Replace it with your
implementation before expecting these checks to pass. Ordinary assembly and
IDE launching resolve the file include beside the caller; the grader's explicit
\texttt{--include} also declares that dependency for its frozen submission.

The original six-function checks remain in the package. Run them against
\texttt{wav-pipeline.asm}, again declaring its helper. Their original caller
check expects filtering without reversal, so it is not the full-pipeline oracle
for the new independent caller. Keep the two intended outputs distinct when
interpreting a failure.

Finally run the two-second generated input through your completed independent
pipeline. Check rate, channels, frame count, complete byte count and exact PCM
against the filtered sequence reversed by whole frames. Compare a fresh run
with a repeat run and inspect the status before opening an output for listening.
Explain why a valid WAV file alone cannot prove that filtering, reversal or
preservation was correct.
